%==============================================================================
%  CAPÍTULO VI
%==============================================================================
\chapter{Procedimiento administrativo de renegociación de la Persona Deudora}
\label{cap:renegociacion}

\fichaCapitulo{Tramitación gratuita ante la \superir{}.}{Ingreso de boletas de
honorarios, audiencias de determinación, renegociación y ejecución.}

%==============================================================================
\bloqueTeorico
%==============================================================================

\subsection{La persona natural como sujeto de consumo financiero}

El procedimiento de renegociación reconoce una asimetría estructural del mercado del
crédito de consumo. La persona natural sobreendeudada no es un empresario que asumió un
riesgo calculado: es, típicamente, un consumidor que contrató crédito en condiciones que
no diseñó ni negoció, frente a instituciones que evaluaron, tarificaron y diversificaron
el riesgo. Cuando el ingreso deja de alcanzar para servir la deuda, la ejecución
individual de cada acreedor produce un resultado socialmente costoso: el deudor queda
capturado en registros negativos que le cierran el acceso al empleo formal y al crédito,
y los propios acreedores recuperan poco de un patrimonio que se disipa en ejecuciones
descoordinadas.

La renegociación ofrece una salida colectiva y ordenada a ese sobreendeudamiento, sin
pasar por el estigma ni el costo de un procedimiento judicial de liquidación.

\subsection{Gratuidad y desjudicialización como reductores del costo social}

Dos rasgos definen el procedimiento y lo distinguen de todo el resto del sistema
concursal. Es \destacar{gratuito} y es \destacar{administrativo}: se tramita ante la
\superir{}, sin tribunal, sin veedor, sin liquidador y sin honorarios de auxiliares.

La decisión de política legislativa es deliberada. Someter el sobreendeudamiento de la
persona natural al procedimiento judicial de liquidación ---con sus costos, su publicidad
y su duración--- habría dejado sin cobertura efectiva a la inmensa mayoría de los deudores
que el sistema pretende atender. La gratuidad y la desjudicialización no son concesiones:
son la condición de acceso que hace que el procedimiento exista para su destinatario real.

\begin{foco}[Qué cambia esto para la asesoría]
El procedimiento está diseñado para tramitarse sin abogado. Ello no significa que la
asistencia letrada carezca de valor ---lo tiene, y considerable, según se desarrolla en el
capítulo \ref{cap:tram-renegociacion}---, sino que el modelo de honorarios debe construirse
sabiendo que el cliente puede prescindir del abogado sin costo. El valor que se cobra es el
del control de admisibilidad, la determinación del pasivo y el diseño de una propuesta
sostenible, no el de una comparecencia que la ley no exige.
\end{foco}

%==============================================================================
\bloqueNormativo
%==============================================================================

\subsection{Causales de admisibilidad}

El \art{260} \parencite{ley-20720} exige la concurrencia copulativa de los siguientes
requisitos:

\begin{enumerate}
  \item Dos o más obligaciones vencidas por más de \destacar{noventa días corridos},
  actualmente exigibles y \destacar{provenientes de obligaciones diversas}.
  \item Que el monto total sea superior a \uf{80}.
  \item Que la persona deudora no haya sido notificada de una demanda que solicite el
  inicio de un procedimiento concursal de liquidación simplificada, ni de cualquier otro
  juicio ejecutivo iniciado en su contra.
\end{enumerate}

\begin{foco}[La excepción laboral en la causal de inadmisibilidad]
El tercer requisito contiene una salvedad que suele pasarse por alto: la exclusión no
opera respecto de los juicios ejecutivos \destacar{de origen laboral}. Una persona
deudora notificada de una ejecución laboral conserva el acceso a la renegociación. La
verificación de causas por RUT debe, por tanto, distinguir la naturaleza de cada
ejecución y no limitarse a constatar su existencia.
\end{foco}

\subsection{Obligaciones excluidas del procedimiento}

El propio \art{260} sustrae determinadas obligaciones del ámbito de la renegociación:

\begin{enumerate}
  \item Los alimentos que se deben por ley, conforme al Título XVIII del Libro I del
  \cc{}, y la compensación económica del Párrafo 1.\textsuperscript{o} del Capítulo VII
  de la \leyn{19.947}{Ley de Matrimonio Civil}.
  \item Las obligaciones derivadas de delitos o cuasidelitos.
\end{enumerate}

\begin{alerta}[Advertencia obligada al cliente]
Estas exclusiones deben comunicarse en la primera reunión y consignarse por escrito. El
deudor que ingresa a la renegociación creyendo que reestructurará su deuda alimentaria
descubrirá lo contrario cuando ya no pueda revertirlo. Véase el capítulo
\ref{cap:entrevista}.
\end{alerta}

\subsection{Ampliación de la legitimación activa}

La \Lreforma{} \citeyearpar{ley-21563} incorporó a las personas naturales que emiten
boletas de honorarios ---contribuyentes de segunda categoría--- corrigiendo una exclusión
histórica que dejaba fuera del sistema a los trabajadores independientes. Sobre el impacto
de la medida en el primer año de vigencia, véase \textcite{superir-balance-primer-ano}.

\subsection{Inicio del procedimiento y solicitud}

El \art{261} dispone que el procedimiento se inicia por la persona deudora ante la
\superir{}, mediante una solicitud cuyo formato está disponible en el sitio web y en las
dependencias del órgano. La solicitud debe acompañar, entre otros antecedentes, una
declaración jurada con la lista de todas las obligaciones del deudor ---vencidas o no,
exigibles o no--- y de todos sus acreedores, con individualización de los montos.

\begin{alerta}[La declaración jurada abarca toda la deuda, no sólo la vencida]
El \art{261} exige declarar la totalidad de las obligaciones, incluidas las no vencidas y
las no exigibles. Omitir deuda futura por creerla irrelevante distorsiona el pasivo y puede
comprometer la seriedad de la solicitud. El barrido documental del capítulo
\ref{cap:entrevista} es aquí el insumo directo de la declaración.
\end{alerta}

\subsection{Examen de admisibilidad}

Conforme al \art{262}, dentro de los \destacar{diez días hábiles administrativos}
siguientes a la presentación, la \superir{} puede: declarar admisible la solicitud, u
ordenar a la persona deudora rectificar sus antecedentes o entregar información adicional,
en el plazo que la propia Superintendencia le fije.

\begin{foco}[Diez días hábiles administrativos, no corridos]
La \Lreforma{} amplió este plazo de cinco a diez días \emph{hábiles administrativos}, que
excluyen sábados, domingos y festivos. El cómputo administrativo difiere del judicial, y la
confusión entre ambos es fuente de error al calcular la respuesta a un requerimiento de
subsanación, cuyo incumplimiento cierra el procedimiento.
\end{foco}

\subsection{Contenido de la Resolución de Admisibilidad}

El \art{263} detalla las menciones de la resolución que declara admisible la solicitud,
entre ellas el nombre y cédula de la persona deudora, el listado inicial de acreedores con
sus montos y preferencias, y el listado de bienes. Esta resolución es la que produce los
efectos protectores del \art{264} y la que fija el pasivo de partida sobre el que operará
la audiencia de determinación.

\subsection{Las tres audiencias del procedimiento}

El procedimiento se estructura en una secuencia de hasta tres audiencias, celebradas ante
el Superintendente o ante quien éste designe.

\begin{description}[leftmargin=0pt,style=nextline,font=\sffamily\bfseries]
  \item[Audiencia de determinación del pasivo (\art{265})] La asistencia es
  \destacar{obligatoria} para todos los acreedores individualizados en la Resolución de
  Admisibilidad y notificados conforme al \art{263}, bajo apercibimiento de proseguir sin
  volver a notificar a los ausentes y de tener por reconocido lo obrado en la audiencia. En
  ella se fija el pasivo sobre el que se negociará.

  \item[Audiencia de renegociación (\art{266})] Determinado el pasivo, se celebra en la
  fecha señalada. El Superintendente facilita la adopción de un acuerdo entre el deudor y
  sus acreedores.

  \item[Audiencia de ejecución (\art{267})] Sólo procede si \destacar{no se alcanzó
  acuerdo} sobre el pasivo o sobre la renegociación. En ella la Superintendencia presenta
  una propuesta de realización del activo del deudor.
\end{description}

\begin{foco}[La asistencia obligatoria del artículo 265 cambia la dinámica]
A diferencia del concurso judicial, donde el acreedor que no comparece simplemente no vota,
aquí el acreedor notificado que no asiste a la determinación del pasivo se expone a que se
tenga por reconocido lo obrado sin él. Es un incentivo poderoso a la comparecencia, y una
herramienta para el deudor cuyo acreedor pretende ignorar el procedimiento.
\end{foco}

\subsection{Efectos de la Resolución de Admisibilidad}

Desde la publicación de la Resolución de Admisibilidad y hasta el término del
procedimiento, el \art{264} produce los siguientes efectos:

\begin{enumerate}
  \item \textbf{Inmunidad ejecutiva} (\Nro 1). No puede solicitarse la liquidación
  forzosa ni voluntaria de la persona deudora, ni iniciarse en su contra juicios
  ejecutivos, ejecuciones de cualquier clase ni restituciones en juicios de arrendamiento.
  Para hacerla valer, el deudor acompaña copia autorizada de la resolución al tribunal
  respectivo, \destacar{y puede comparecer personalmente, sin patrocinio de abogado}. La
  ley precisa que sólo puede hacerse valer como excepción.

  \item \textbf{Suspensión de la prescripción} (\Nro 2). Se suspenden los plazos de
  prescripción extintiva de las obligaciones del deudor.

  \item \textbf{Cese del devengo de intereses moratorios} (\Nro 3). No continúan
  devengándose los intereses moratorios pactados en los actos y contratos vigentes.
\end{enumerate}

\begin{practica}[Un efecto de valor inmediato para el cliente]
El cese del devengo de intereses moratorios detiene el crecimiento del pasivo desde la
admisibilidad. En deudas de consumo con mora prolongada, este solo efecto puede
representar una porción sustancial del beneficio del procedimiento, y conviene explicarlo
al cliente en términos cuantitativos.
\end{practica}

\subsection{El acuerdo de ejecución, el plan de reembolso y sus quórums}

En la audiencia de ejecución del \art{267}, la persona deudora y \destacar{dos o más
acreedores que representen al menos el 50\,\% del pasivo reconocido con derecho a voto}
acuerdan la fórmula de realización del activo. No se computan, para el quórum ni para las
votaciones, los créditos de las personas relacionadas con el deudor.

El acuerdo puede contener un plan de reembolso sujeto a dos límites: no puede exceder de
\destacar{seis meses} de duración, y la cuota mensual no puede ser superior al
\destacar{30\,\% de los ingresos del deudor}. La \superir{} puede suspender la audiencia
hasta por diez días.

\begin{alerta}[Corrección de una atribución errónea]
El plan de reembolso y la fórmula de realización se regulan en el \art{267} ---audiencia de
ejecución--- y no en el \art{264}, que trata de los efectos de la resolución de
admisibilidad. La confusión es frecuente en la literatura de divulgación.
\end{alerta}

\subsection{Finalización del procedimiento y extinción de saldos}

El \art{268} regula la resolución que declara finalizado el procedimiento, una vez vencido
el plazo para impugnar el acuerdo ---de renegociación o de ejecución--- o desechada la
impugnación conforme al \art{272}.

El efecto es de la mayor importancia y conviene subrayarlo: si el procedimiento finaliza en
virtud de un \destacar{acuerdo de ejecución}, se entienden extinguidos por el solo
ministerio de la ley los saldos insolutos de las obligaciones. Es el equivalente
administrativo del \discharge{} judicial (capítulo \ref{cap:discharge-cae}).

\begin{foco}[La renegociación también puede extinguir deuda]
Suele presentarse la renegociación como un mero reordenamiento de plazos. No lo es: cuando
desemboca en un acuerdo de ejecución, extingue los saldos insolutos igual que una
liquidación. La diferencia con la liquidación simplificada está en el camino ---negociado y
administrativo, no forzoso y judicial---, no necesariamente en el resultado extintivo.
\end{foco}

\subsection{El término anticipado del procedimiento}

El \art{269} faculta a la \superir{} para declarar el \destacar{término anticipado} de la
renegociación en dos supuestos principales: si la persona deudora infringe la prohibición de
disponer de sus bienes del \art{264} \Nro 6 ---sin perjuicio de la sanción propia del depositario
alzado del artículo 444 del \cpc{}---, o si deja de cumplir los requisitos de admisibilidad.

\begin{alerta}[Disponer de los bienes durante el procedimiento cierra la puerta]
Admitida la renegociación, el deudor no puede disponer libremente de sus bienes. Enajenar un
activo declarado durante el procedimiento no sólo provoca el término anticipado: activa la
sanción del depositario alzado. Es una advertencia que debe hacerse expresamente al cliente, que
suele ignorar que la admisibilidad le impuso esa restricción.
\end{alerta}

\subsection{La impugnación del acuerdo}

El \art{272} permite a los acreedores afectados impugnar el acuerdo de renegociación o de
ejecución, por causales tasadas: \destacar{error en el cómputo de las mayorías} que incida en el
quórum, o \destacar{falsedad o exageración del crédito} de alguno de los acreedores que concurrió
a formarlo. Es el equivalente administrativo de la impugnación del acuerdo de reorganización
(capítulo \ref{cap:reorganizacion-ordinaria}, \art{85}).

\subsection{El régimen de recursos administrativos}

Por tratarse de un procedimiento administrativo, el régimen recursivo del \art{270} es distinto
del judicial: contra la resolución que finaliza el procedimiento ---o que lo termina
anticipadamente--- procede el \destacar{recurso de reposición administrativa} del artículo 59 de
la \leyn{19.880}{Ley de Bases de los Procedimientos Administrativos}, y contra la resolución que
lo desecha, el \destacar{recurso de reclamación} del \art{341}. La interposición de estos
recursos no suspende, por regla general, los efectos de lo resuelto.

\begin{foco}[No confundir la vía administrativa con la judicial]
El error frecuente es intentar apelar ante un tribunal una decisión de la \superir{}. La vía es
administrativa: reposición ante el propio órgano y luego reclamación. Conocer los plazos y la
secuencia del artículo 59 de la Ley 19.880 ---distintos de los procesales--- es indispensable para
no dejar firme una resolución adversa.
\end{foco}

\subsection{Bienes inembargables en el acuerdo de ejecución}

El \art{271} preserva la inembargabilidad de los bienes del artículo 445 del \cpc{} y de los que las leyes
declaren inembargables. Si la persona deudora es casada, la realización de sus bienes se sujeta a
las normas del \cc{} según el régimen matrimonial pactado.

\subsection{El procedimiento en cifras: el impacto empírico de la reforma}

La ampliación del acceso a las personas que emiten boletas de honorarios convirtió a la
renegociación en el procedimiento de mayor crecimiento del sistema. Las estadísticas de la
\superir{} \parencite{superir-estadisticas-2025} muestran la magnitud del fenómeno.

Entre enero y septiembre de 2025 se iniciaron en Chile \destacar{8.475 procedimientos
concursales}, distribuidos como sigue:

\begin{table}[htbp]
\centering
\small
\caption{Procedimientos concursales iniciados (enero--septiembre de 2025)}
\label{tab:cap06-estadisticas}
\begin{tabularx}{\textwidth}{@{}>{\raggedright\arraybackslash}X p{0.16\textwidth} p{0.16\textwidth}@{}}
\toprule
\textbf{Tipo de procedimiento} & \textbf{Casos} & \textbf{\% del total} \\
\midrule
Liquidación simplificada de persona deudora & 5.024 & 59{,}3\,\% \\
Renegociación de persona deudora & 2.891 & 34{,}1\,\% \\
Liquidación simplificada MIPE & 387 & 4{,}6\,\% \\
Liquidación ordinaria de mediana y gran empresa & 117 & 1{,}4\,\% \\
Liquidación ordinaria de persona deudora & 24 & 0{,}3\,\% \\
Reorganización judicial de mediana y gran empresa & 24 & 0{,}3\,\% \\
Reorganización simplificada MIPE & 8 & 0{,}1\,\% \\
\midrule
\textbf{Total} & \textbf{8.475} & \textbf{100\,\%} \\
\bottomrule
\end{tabularx}
\end{table}

La lectura de la tabla es contundente: \destacar{más del 93\,\% del sistema concursal chileno son
procedimientos de persona deudora} ---liquidación simplificada y renegociación---, mientras que la
reorganización de empresas, ordinaria y simplificada, apenas reúne 32 casos en nueve meses. El
derecho concursal chileno es, en los hechos, un derecho de la insolvencia de la persona natural.

Un segundo grupo de indicadores describe la \destacar{efectividad de las audiencias} de la
renegociación: se alcanza acuerdo en el 99{,}4\,\% de las audiencias de determinación del pasivo y
en el 98{,}5\,\% de las de renegociación, mientras que la audiencia de ejecución ---que sólo opera
cuando la renegociación fracasa--- registra un 6{,}7\,\%. La inmensa mayoría de los procedimientos
se cierra por acuerdo en las dos primeras audiencias, y la ejecución queda como válvula residual.

En cuanto al perfil de los concursados, las mujeres representan el \destacar{42{,}0\,\%} de las
liquidaciones simplificadas de personas naturales. En el plano territorial, la \destacar{Región
Metropolitana} concentra el 66{,}7\,\% de las liquidaciones ordinarias de personas y el 46{,}9\,\%
de las renegociaciones, seguida por Biobío (12{,}3\,\%) y Valparaíso (8{,}4\,\%): una concentración
que refleja tanto la densidad poblacional como la brecha de cobertura del procedimiento en
regiones.\verificar{Actualizar todas estas cifras contra el boletín estadístico oficial de la
\superir{} a la fecha de cierre de la edición.}

\begin{foco}[Qué dicen los números sobre el procedimiento]
Cuatro lecturas. Primera: el sistema concursal chileno es abrumadoramente un sistema de persona
deudora ---más del 93\,\% de los casos---, y la reorganización empresarial sigue siendo marginal.
Segunda: la efectividad de la renegociación es altísima (sobre 98\,\% de acuerdos en las dos
primeras audiencias), lo que confirma que cumple su función de evitar la liquidación. Tercera: la
participación femenina, cercana a la mitad, muestra que el procedimiento alcanza a un segmento que
el sistema anterior no atendía. Cuarta: la concentración metropolitana revela una brecha de acceso
regional que la política pública aún no cierra. Fue este volumen el que motivó la dictación de la
\ncg{28} \citeyearpar{ncg-28} para precisar y agilizar los criterios de admisibilidad.
\end{foco}

\subsection{Espejo comparado: la reestructuración de deudas de la persona natural}

La renegociación chilena tiene equivalentes funcionales en los grandes ordenamientos, cuyo
contraste ilumina sus opciones de diseño.

\subsubsection{El \emph{Chapter 13} estadounidense}

El \emph{Chapter 13} del \emph{Bankruptcy Code} \parencite{us-bankruptcy-code} ---«ajuste de las
deudas de un individuo con ingreso regular»--- permite a la persona natural con ingreso estable
\destacar{conservar sus bienes} y reestructurar sus deudas mediante un \destacar{plan de pagos de
tres a cinco años} con cargo a su ingreso futuro, administrado por un \emph{trustee}, con
\discharge{} al completarlo. Es la alternativa a la liquidación del \emph{Chapter 7}.

El paralelo con Chile es imperfecto pero instructivo: el plan de reembolso del \art{267} cumple
una función análoga ---pago ordenado con cargo al ingreso--- pero es mucho más breve (seis meses,
no cinco años) y opera en sede administrativa, no judicial.

\subsubsection{La «segunda oportunidad» española}

Tras la \leyn{16/2022}{reforma española del texto refundido concursal}
\parencite{esp-ley16-2022}, España ofrece a la persona natural dos vías de exoneración del pasivo
insatisfecho: la \destacar{exoneración inmediata} por liquidación de la masa, o la sujeción a un
\destacar{plan de pagos de tres a cinco años} que permite \destacar{conservar la vivienda
habitual}. El acceso exige un estándar de buena fe del deudor.

\begin{foco}[Chile eligió la vía rápida y administrativa]
Frente al plan de pagos plurianual de EE.UU. y España, Chile optó por un procedimiento
\emph{breve} (seis meses) y \emph{administrativo} (ante la \superir{}, sin tribunal). La ventaja
es la celeridad y el bajo costo; el límite, que un plan de seis meses difícilmente resuelve un
sobreendeudamiento profundo, que en los modelos comparados se distribuye en tres a cinco años.
La discusión de política pública ---plazo corto y masivo versus plan largo y personalizado---
está abierta.
\end{foco}

%==============================================================================
\bloquePractico
%==============================================================================

\subsection{Tramitación en la plataforma de la \superir{}}

\begin{flujograma}{Secuencia del procedimiento administrativo de renegociación}{cap06-iter}
  \node[hito] (ingreso) {SOLICITUD ANTE LA \superir{} (\art{261})\\[2pt]
    \normalfont Declaración jurada de todas las obligaciones};
  \node[nodo, below=of ingreso] (examen)
    {EXAMEN DE ADMISIBILIDAD (\art{262})\\[2pt]
     \normalfont 10 días hábiles administrativos};
  \node[nodo, below=of examen] (resol)
    {RESOLUCIÓN DE ADMISIBILIDAD (\art{263})\\[2pt]
     \normalfont Efectos del \art{264}: suspensión de ejecuciones,
     prescripción e intereses moratorios};
  \node[nodo, below=of resol] (pasivo)
    {AUDIENCIA DE DETERMINACIÓN DEL PASIVO (\art{265})\\[2pt]
     \normalfont Asistencia obligatoria de los acreedores};
  \node[hito, below=of pasivo] (reneg)
    {AUDIENCIA DE RENEGOCIACIÓN (\art{266})};

  \coordinate (split) at ($(reneg.south)+(0,-6mm)$);
  \node[favorable, text width=40mm, below left=12mm and 2mm of reneg.south] (aprueba)
    {ACUERDO DE RENEGOCIACIÓN};
  \node[adverso, text width=40mm, below right=12mm and 2mm of reneg.south] (ejec)
    {AUDIENCIA DE EJECUCIÓN (\art{267})\\[2pt]
     \normalfont Propuesta de realización; quórum $\geq$50\,\%};

  \node[favorable, text width=40mm, below=9mm of aprueba] (fin1)
    {FINALIZACIÓN (\art{268})};
  \node[favorable, text width=40mm, below=9mm of ejec] (fin2)
    {ACUERDO DE EJECUCIÓN\\[2pt]
     \normalfont Extingue los saldos insolutos (\art{268})};

  \draw[flecha] (ingreso) -- (examen);
  \draw[flecha] (examen) -- (resol);
  \draw[flecha] (resol) -- (pasivo);
  \draw[flecha] (pasivo) -- (reneg);
  \draw[flecha] (reneg.south) -- (split) -| (aprueba.north);
  \draw[flecha] (reneg.south) -- (split) -| (ejec.north);
  \draw[flecha] (aprueba) -- (fin1);
  \draw[flecha] (ejec) -- (fin2);
\end{flujograma}

\subsection{Calendario de plazos}

\begin{table}[htbp]
\centering
\small
\caption{Plazos del procedimiento de renegociación}
\label{tab:plazos-renegociacion}
\begin{tabularx}{\textwidth}{@{}>{\raggedright\arraybackslash}p{0.40\textwidth} p{0.22\textwidth} X@{}}
\toprule
\textbf{Actuación} & \textbf{Plazo} & \textbf{Norma} \\
\midrule
Examen de admisibilidad & 10 días hábiles administrativos & \art{262} \\
\addlinespace
Subsanación de antecedentes & El que fije la \superir{} & \art{262} \Nro 2 \\
\addlinespace
Plan de reembolso (duración) & Máximo 6 meses & \art{267} \\
\addlinespace
Plan de reembolso (cuota) & Máximo 30\,\% de los ingresos & \art{267} \\
\addlinespace
Suspensión de audiencia & Hasta 10 días & \art{267} \\
\bottomrule
\end{tabularx}
\end{table}

\begin{alerta}[Plazos hábiles administrativos]
Los plazos de este procedimiento se cuentan en días \emph{hábiles administrativos}
(excluidos sábados, domingos y festivos), y no en días corridos ni en el cómputo judicial.
Es una diferencia que altera de forma relevante el calendario real.
\end{alerta}

\subsection{Preparación de las audiencias}

Remisión al capítulo \ref{cap:tram-renegociacion}, que desarrolla la construcción de la
nómina propia, las objeciones al pasivo en la audiencia de determinación y el diseño de una
propuesta sostenible sobre la capacidad real de pago.

%==============================================================================
\bloqueJurisprudencial
%==============================================================================

\subsection{Doctrina administrativa de la \superir{}}

Pautas de la \ncg{28} \citeyearpar{ncg-28} sobre acreditación de la existencia y monto de las obligaciones de
codeudores solidarios y avales dentro del proceso de renegociación.

\subsection{Recursos de protección por mantención de registros negativos}

Fallos que sancionan a instituciones financieras que mantienen vigentes registros
comerciales negativos respecto de deudores con acuerdos de renegociación firmes.
\verificar{Individualizar sentencias.}
